\begin{PSbox}[unbreakable]{obj}{Learning Objectives}

After completing this module, students will be able to:

\begin{enumerate}
\def\labelenumi{\arabic{enumi}.}
\tightlist
\item
  Apply linear transformations to random variables
\item
  Derive PDFs of nonlinear transformations using the CDF and Jacobian methods
\item
  Apply transformations to engineering problems (power, SNR, envelope)
\item
  Verify analytical results with MATLAB simulation
\end{enumerate}

\end{PSbox}

\PSsection{5.1}{Introduction: Why Transform Random Variables?}

In engineering, we often know the statistics of one quantity but need the statistics of a \textbf{derived} quantity:

{\def\LTcaptype{none} % do not increment counter
\begin{longtable}[c]{@{}ccc@{}}
\toprule\noalign{}
\rowcolor{PSnavy}\PSth{Known} & \PSth{Needed} & \PSth{Transformation} \\
\midrule\noalign{}
\endhead
\bottomrule\noalign{}
\endlastfoot
Voltage V & Power P & P = V\textsuperscript{2}/R \\
Linear SNR & SNR in dB & Y = 10\ensuremath{\cdot}log\textsubscript{10}(X) \\
Gaussian I, Q & Envelope & R = \ensuremath{\surd}(I\textsuperscript{2} + Q\textsuperscript{2}) \\
Current I & Power & P = I\textsuperscript{2}R \\
Distance d & Path loss & L = d\textsuperscript{\ensuremath{\alpha}} \\
\end{longtable}
}

\PSsection{5.2}{Linear Transformations: Y = aX + b}

\PSsubsection{Result}

If X has PDF f\textsubscript{X}(x), then Y = aX + b has:

\PSdisplay{f_Y(y) = \frac{1}{|a|} f_X\left(\frac{y-b}{a}\right)}

\PSsubsection{Mean and Variance}

\begin{itemize}
\tightlist
\item
  E[Y] = aE[X] + b
\item
  Var(Y) = a\textsuperscript{2}Var(X)
\end{itemize}

\begin{PSbox}[unbreakable]{ex}{Engineering Example: Amplifier}

Input signal X \textasciitilde{} N(0, \ensuremath{\sigma}\textsuperscript{2}) passes through amplifier with gain G = 10 and DC offset V\textsubscript{0} = 1.5V.

Output: Y = 10X + 1.5

\begin{itemize}
\tightlist
\item
  Y \textasciitilde{} N(10\ensuremath{\cdot}0 + 1.5, 10\textsuperscript{2}\ensuremath{\cdot}\ensuremath{\sigma}\textsuperscript{2}) = N(1.5, 100\ensuremath{\sigma}\textsuperscript{2})
\item
  Output noise power increased by G\textsuperscript{2} = 100
\end{itemize}

\end{PSbox}

\begin{PSbox}[unbreakable]{ex}{Engineering Example: Temperature Sensor}

Sensor output V = 0.01\ensuremath{\cdot}T + 0.5 (V in volts, T in \ensuremath{^{\circ}}C)

If T \textasciitilde{} N(25, 4): V \textasciitilde{} N(0.01\ensuremath{\times}25 + 0.5, 0.01\textsuperscript{2}\ensuremath{\times}4) = N(0.75, 0.0004)

\end{PSbox}

\PSsection{5.3}{General Monotonic Transformations}

\PSsubsection{CDF Method (Most General)}

For Y = g(X), find f\textsubscript{Y}(y) by:

\begin{enumerate}
\def\labelenumi{\arabic{enumi}.}
\tightlist
\item
  Write F\textsubscript{Y}(y) = P(Y \ensuremath{\le} y) = P(g(X) \ensuremath{\le} y)
\item
  Solve the inequality for X
\item
  Express in terms of F\textsubscript{X}
\item
  Differentiate to get f\textsubscript{Y}(y)
\end{enumerate}

\PSsubsection{Formula for Monotonic g}

If g is strictly increasing with inverse g\textsuperscript{-1}:\PSdisplay{f_Y(y) = f_X(g^{-1}(y)) \cdot \frac{d}{dy}[g^{-1}(y)]}

If g is strictly decreasing:\PSdisplay{f_Y(y) = f_X(g^{-1}(y)) \cdot \left|\frac{d}{dy}[g^{-1}(y)]\right|}

Combined (works for both):\PSdisplay{f_Y(y) = \frac{f_X(x)}{|g'(x)|}\bigg|_{x=g^{-1}(y)}}

\begin{PSbox}[unbreakable]{ex}{Engineering Example: SNR in dB}

Linear SNR: X \textasciitilde{} Exponential(1) (normalized Rayleigh fading power).\PSbr{}dB SNR: Y = 10\ensuremath{\cdot}log\textsubscript{10}(X)

\begin{itemize}
\tightlist
\item
  g(x) = 10\ensuremath{\cdot}log\textsubscript{10}(x), g’(x) = 10/(x\ensuremath{\cdot}ln10)
\item
  g\textsuperscript{-1}(y) = 10\textsuperscript{y/10}
\end{itemize}

\PSdisplay{f_Y(y) = \frac{f_X(10^{y/10})}{10/(10^{y/10} \cdot \ln 10)} = \frac{\ln 10}{10} \cdot 10^{y/10} \cdot e^{-10^{y/10}}}

This is NOT Gaussian — it has a longer left tail (deep fades).

\end{PSbox}

\PSsection{5.4}{Nonlinear Transformations: Non-Monotonic Case}

\PSsubsection{When Y = g(X) is Not One-to-One}

If multiple x-values map to the same y, sum over all solutions:

\PSdisplay{f_Y(y) = \sum_{i} \frac{f_X(x_i)}{|g'(x_i)|}}

where x\textsubscript{1}, x\textsubscript{2}, … are all roots of g(x) = y.

\begin{PSbox}[unbreakable]{ex}{Engineering Example: Power from Voltage (Y = X\textsuperscript{2})}

Voltage X \textasciitilde{} N(0, \ensuremath{\sigma}\textsuperscript{2}). Power: Y = X\textsuperscript{2}.

For y \textgreater{} 0, solutions are x = +\ensuremath{\surd}y and x = -\ensuremath{\surd}y. g’(x) = 2x.

\PSdisplay{f_Y(y) = \frac{f_X(\sqrt{y})}{2\sqrt{y}} + \frac{f_X(-\sqrt{y})}{2\sqrt{y}} = \frac{1}{\sqrt{2\pi}\sigma} \cdot \frac{e^{-y/(2\sigma^2)}}{\sqrt{y}}}

This is a \textbf{Chi-squared distribution with 1 degree of freedom} (scaled by \ensuremath{\sigma}\textsuperscript{2}).

\end{PSbox}

\begin{PSbox}[unbreakable]{ex}{Engineering Example: Full-Wave Rectifier (Y = \textbar{}X\textbar{})}

Input X \textasciitilde{} N(0, \ensuremath{\sigma}\textsuperscript{2}). Output Y = \textbar{}X\textbar{}.

For y \textgreater{} 0: x = +y and x = -y are solutions. \textbar{}g’(x)\textbar{} = 1.

\PSdisplay{f_Y(y) = f_X(y) + f_X(-y) = \frac{2}{\sigma\sqrt{2\pi}} e^{-y^2/(2\sigma^2)}, \quad y \geq 0}

This is the \textbf{folded normal} (half-normal) distribution.

\end{PSbox}

\PSsection{5.5}{The Jacobian Method}

\PSsubsection{Procedure for Y = g(X)}

\begin{enumerate}
\def\labelenumi{\arabic{enumi}.}
\tightlist
\item
  Identify the transformation y = g(x)
\item
  Find the inverse: x = g\textsuperscript{-1}(y)
\item
  Compute the Jacobian: J = \textbar{}dx/dy\textbar{} = \textbar{}d[g\textsuperscript{-1}(y)]/dy\textbar{}
\item
  Apply: f\textsubscript{Y}(y) = f\textsubscript{X}(g\textsuperscript{-1}(y)) \ensuremath{\cdot} \textbar{}J\textbar{}
\end{enumerate}

\PSsubsection{Step-by-Step Example: Exponential of Gaussian}

X \textasciitilde{} N(\ensuremath{\mu}, \ensuremath{\sigma}\textsuperscript{2}). Y = e\textsuperscript{x} (log-normal transformation).

\begin{enumerate}
\def\labelenumi{\arabic{enumi}.}
\tightlist
\item
  g(x) = e\textsuperscript{x} (monotonically increasing)
\item
  x = ln(y), valid for y \textgreater{} 0
\item
  J = \textbar{}dx/dy\textbar{} = 1/y
\item
  f\textsubscript{Y}(y) = f\textsubscript{X}(ln y) \ensuremath{\cdot} (1/y) = (1/(y\ensuremath{\sigma}\ensuremath{\surd}(2\ensuremath{\pi}))) \ensuremath{\cdot} exp(-(ln y - \ensuremath{\mu})\textsuperscript{2}/(2\ensuremath{\sigma}\textsuperscript{2}))
\end{enumerate}

This is the \textbf{log-normal distribution} — models many engineering quantities (shadowing in wireless, stock prices, component lifetimes with wear).

\PSsection{5.6}{Engineering Application: Rayleigh from Gaussian}

\PSsubsection{Derivation of Rayleigh Distribution}

In communications, the received signal has in-phase (I) and quadrature (Q) components:

\begin{itemize}
\tightlist
\item
  I \textasciitilde{} N(0, \ensuremath{\sigma}\textsuperscript{2}), Q \textasciitilde{} N(0, \ensuremath{\sigma}\textsuperscript{2}), independent
\end{itemize}

Envelope: R = \ensuremath{\surd}(I\textsuperscript{2} + Q\textsuperscript{2})

Using the transformation from (I, Q) \ensuremath{\to} (R, \ensuremath{\theta}) with polar coordinates:

\begin{itemize}
\tightlist
\item
  I = R\ensuremath{\cdot}cos(\ensuremath{\theta}), Q = R\ensuremath{\cdot}sin(\ensuremath{\theta})
\item
  Jacobian: \textbar{}\ensuremath{\partial}(I,Q)/\ensuremath{\partial}(R,\ensuremath{\theta})\textbar{} = R
\end{itemize}

Joint PDF of (R, \ensuremath{\theta}):\PSbr{}f\textsubscript{R,\ensuremath{\Theta}}(r,\ensuremath{\theta}) = f\textsubscript{I,Q}(r\ensuremath{\cdot}cos\ensuremath{\theta}, r\ensuremath{\cdot}sin\ensuremath{\theta}) \ensuremath{\cdot} r = (r/(2\ensuremath{\pi}\ensuremath{\sigma}\textsuperscript{2})) \ensuremath{\cdot} e\textsuperscript{-r\textsuperscript{2}/(2\ensuremath{\sigma}\textsuperscript{2})}

Marginalizing over \ensuremath{\theta} \ensuremath{\in} [0, 2\ensuremath{\pi}):

\PSdisplay{f_R(r) = \frac{r}{\sigma^2} e^{-r^2/(2\sigma^2)}, \quad r \geq 0}

This is the \textbf{Rayleigh distribution} — fundamental to wireless fading channel modeling.

\PSsection{5.7}{MATLAB Examples}

\begin{PSbox}[unbreakable]{ex}{Example 1: Linear Transformation Verification}

\tcbset{PScodemode/.style={unbreakable}}

\begin{Shaded}
\begin{Highlighting}[]
\CommentTok{\%\% Linear Transformation: Amplifier Output}
\VariableTok{sigma\_in} \OperatorTok{=} \FloatTok{0.1}\OperatorTok{;}  \VariableTok{G} \OperatorTok{=} \FloatTok{5}\OperatorTok{;}  \VariableTok{offset} \OperatorTok{=} \FloatTok{2}\OperatorTok{;}
\VariableTok{N} \OperatorTok{=} \FloatTok{100000}\OperatorTok{;}

\VariableTok{X} \OperatorTok{=} \VariableTok{sigma\_in} \OperatorTok{*} \VariableTok{randn}\NormalTok{(}\FloatTok{1}\OperatorTok{,} \VariableTok{N}\NormalTok{)}\OperatorTok{;}   \CommentTok{\% Input noise}
\VariableTok{Y} \OperatorTok{=} \VariableTok{G}\OperatorTok{*}\VariableTok{X} \OperatorTok{+} \VariableTok{offset}\OperatorTok{;}              \CommentTok{\% Output}

\VariableTok{figure}\OperatorTok{;}
\VariableTok{subplot}\NormalTok{(}\FloatTok{1}\OperatorTok{,}\FloatTok{2}\OperatorTok{,}\FloatTok{1}\NormalTok{)}\OperatorTok{;}
\VariableTok{histogram}\NormalTok{(}\VariableTok{X}\OperatorTok{,} \FloatTok{80}\OperatorTok{,} \SpecialStringTok{\textquotesingle{}Normalization\textquotesingle{}}\OperatorTok{,} \SpecialStringTok{\textquotesingle{}pdf\textquotesingle{}}\NormalTok{)}\OperatorTok{;} \VariableTok{hold} \VariableTok{on}\OperatorTok{;}
\VariableTok{x} \OperatorTok{=} \VariableTok{linspace}\NormalTok{(}\OperatorTok{{-}}\FloatTok{0.5}\OperatorTok{,} \FloatTok{0.5}\OperatorTok{,} \FloatTok{200}\NormalTok{)}\OperatorTok{;}
\VariableTok{plot}\NormalTok{(}\VariableTok{x}\OperatorTok{,} \VariableTok{normpdf}\NormalTok{(}\VariableTok{x}\OperatorTok{,} \FloatTok{0}\OperatorTok{,} \VariableTok{sigma\_in}\NormalTok{)}\OperatorTok{,} \SpecialStringTok{\textquotesingle{}r\textquotesingle{}}\OperatorTok{,} \SpecialStringTok{\textquotesingle{}LineWidth\textquotesingle{}}\OperatorTok{,} \FloatTok{2}\NormalTok{)}\OperatorTok{;}
\VariableTok{title}\NormalTok{(}\SpecialStringTok{\textquotesingle{}Input X \textasciitilde{} N(0, 0.01)\textquotesingle{}}\NormalTok{)}\OperatorTok{;} \VariableTok{xlabel}\NormalTok{(}\SpecialStringTok{\textquotesingle{}Voltage (V)\textquotesingle{}}\NormalTok{)}\OperatorTok{;}

\VariableTok{subplot}\NormalTok{(}\FloatTok{1}\OperatorTok{,}\FloatTok{2}\OperatorTok{,}\FloatTok{2}\NormalTok{)}\OperatorTok{;}
\VariableTok{histogram}\NormalTok{(}\VariableTok{Y}\OperatorTok{,} \FloatTok{80}\OperatorTok{,} \SpecialStringTok{\textquotesingle{}Normalization\textquotesingle{}}\OperatorTok{,} \SpecialStringTok{\textquotesingle{}pdf\textquotesingle{}}\NormalTok{)}\OperatorTok{;} \VariableTok{hold} \VariableTok{on}\OperatorTok{;}
\VariableTok{y} \OperatorTok{=} \VariableTok{linspace}\NormalTok{(}\FloatTok{0}\OperatorTok{,} \FloatTok{4}\OperatorTok{,} \FloatTok{200}\NormalTok{)}\OperatorTok{;}
\VariableTok{plot}\NormalTok{(}\VariableTok{y}\OperatorTok{,} \VariableTok{normpdf}\NormalTok{(}\VariableTok{y}\OperatorTok{,} \VariableTok{offset}\OperatorTok{,} \VariableTok{G}\OperatorTok{*}\VariableTok{sigma\_in}\NormalTok{)}\OperatorTok{,} \SpecialStringTok{\textquotesingle{}r\textquotesingle{}}\OperatorTok{,} \SpecialStringTok{\textquotesingle{}LineWidth\textquotesingle{}}\OperatorTok{,} \FloatTok{2}\NormalTok{)}\OperatorTok{;}
\VariableTok{title}\NormalTok{(}\SpecialStringTok{\textquotesingle{}Output Y = 5X + 2\textquotesingle{}}\NormalTok{)}\OperatorTok{;} \VariableTok{xlabel}\NormalTok{(}\SpecialStringTok{\textquotesingle{}Voltage (V)\textquotesingle{}}\NormalTok{)}\OperatorTok{;}
\end{Highlighting}
\end{Shaded}

\end{PSbox}

\begin{PSbox}[unbreakable]{ex}{Example 2: Power from Gaussian Voltage}

\tcbset{PScodemode/.style={unbreakable}}

\begin{Shaded}
\begin{Highlighting}[]
\CommentTok{\%\% Nonlinear: Y = X\^{}2 (instantaneous power)}
\VariableTok{sigma} \OperatorTok{=} \FloatTok{1}\OperatorTok{;} \VariableTok{N} \OperatorTok{=} \FloatTok{200000}\OperatorTok{;}
\VariableTok{X} \OperatorTok{=} \VariableTok{sigma} \OperatorTok{*} \VariableTok{randn}\NormalTok{(}\FloatTok{1}\OperatorTok{,} \VariableTok{N}\NormalTok{)}\OperatorTok{;}
\VariableTok{Y} \OperatorTok{=} \VariableTok{X}\OperatorTok{.\^{}}\FloatTok{2}\OperatorTok{;}

\VariableTok{figure}\OperatorTok{;}
\VariableTok{histogram}\NormalTok{(}\VariableTok{Y}\OperatorTok{,} \FloatTok{150}\OperatorTok{,} \SpecialStringTok{\textquotesingle{}Normalization\textquotesingle{}}\OperatorTok{,} \SpecialStringTok{\textquotesingle{}pdf\textquotesingle{}}\OperatorTok{,} \SpecialStringTok{\textquotesingle{}BinLimits\textquotesingle{}}\OperatorTok{,}\NormalTok{ [}\FloatTok{0}\OperatorTok{,} \FloatTok{8}\NormalTok{])}\OperatorTok{;}
\VariableTok{hold} \VariableTok{on}\OperatorTok{;}
\VariableTok{y} \OperatorTok{=} \VariableTok{linspace}\NormalTok{(}\FloatTok{0.01}\OperatorTok{,} \FloatTok{8}\OperatorTok{,} \FloatTok{300}\NormalTok{)}\OperatorTok{;}
\VariableTok{f\_theory} \OperatorTok{=} \VariableTok{chi2pdf}\NormalTok{(}\VariableTok{y}\OperatorTok{/}\VariableTok{sigma}\OperatorTok{\^{}}\FloatTok{2}\OperatorTok{,} \FloatTok{1}\NormalTok{) }\OperatorTok{/} \VariableTok{sigma}\OperatorTok{\^{}}\FloatTok{2}\OperatorTok{;}
\VariableTok{plot}\NormalTok{(}\VariableTok{y}\OperatorTok{,} \VariableTok{f\_theory}\OperatorTok{,} \SpecialStringTok{\textquotesingle{}r\textquotesingle{}}\OperatorTok{,} \SpecialStringTok{\textquotesingle{}LineWidth\textquotesingle{}}\OperatorTok{,} \FloatTok{2}\NormalTok{)}\OperatorTok{;}
\VariableTok{xlabel}\NormalTok{(}\SpecialStringTok{\textquotesingle{}Power (V²)\textquotesingle{}}\NormalTok{)}\OperatorTok{;} \VariableTok{ylabel}\NormalTok{(}\SpecialStringTok{\textquotesingle{}PDF\textquotesingle{}}\NormalTok{)}\OperatorTok{;}
\VariableTok{title}\NormalTok{(}\SpecialStringTok{\textquotesingle{}Distribution of Power Y = X²\textquotesingle{}}\NormalTok{)}\OperatorTok{;}
\VariableTok{legend}\NormalTok{(}\SpecialStringTok{\textquotesingle{}Simulation\textquotesingle{}}\OperatorTok{,} \SpecialStringTok{\textquotesingle{}Theory (scaled χ²₁)\textquotesingle{}}\NormalTok{)}\OperatorTok{;}
\end{Highlighting}
\end{Shaded}

\end{PSbox}

\begin{PSbox}[unbreakable]{ex}{Example 3: SNR in dB from Linear SNR}

\tcbset{PScodemode/.style={unbreakable}}

\begin{Shaded}
\begin{Highlighting}[]
\CommentTok{\%\% SNR transformation: Y = 10*log10(X)}
\CommentTok{\% X \textasciitilde{} Exponential(1) (Rayleigh fading power, normalized)}
\VariableTok{N} \OperatorTok{=} \FloatTok{200000}\OperatorTok{;}
\VariableTok{X} \OperatorTok{=} \VariableTok{exprnd}\NormalTok{(}\FloatTok{1}\OperatorTok{,} \FloatTok{1}\OperatorTok{,} \VariableTok{N}\NormalTok{)}\OperatorTok{;}
\VariableTok{Y} \OperatorTok{=} \FloatTok{10}\OperatorTok{*}\VariableTok{log10}\NormalTok{(}\VariableTok{X}\NormalTok{)}\OperatorTok{;}   \CommentTok{\% SNR in dB}

\VariableTok{figure}\OperatorTok{;}
\VariableTok{histogram}\NormalTok{(}\VariableTok{Y}\OperatorTok{,} \FloatTok{150}\OperatorTok{,} \SpecialStringTok{\textquotesingle{}Normalization\textquotesingle{}}\OperatorTok{,} \SpecialStringTok{\textquotesingle{}pdf\textquotesingle{}}\NormalTok{)}\OperatorTok{;}
\VariableTok{hold} \VariableTok{on}\OperatorTok{;}
\VariableTok{y} \OperatorTok{=} \VariableTok{linspace}\NormalTok{(}\OperatorTok{{-}}\FloatTok{30}\OperatorTok{,} \FloatTok{15}\OperatorTok{,} \FloatTok{300}\NormalTok{)}\OperatorTok{;}
\CommentTok{\% Theoretical PDF of Y = 10*log10(X) where X \textasciitilde{} Exp(1)}
\VariableTok{f\_Y} \OperatorTok{=}\NormalTok{ (}\VariableTok{log}\NormalTok{(}\FloatTok{10}\NormalTok{)}\OperatorTok{/}\FloatTok{10}\NormalTok{) }\OperatorTok{.*} \FloatTok{10}\OperatorTok{.\^{}}\NormalTok{(}\VariableTok{y}\OperatorTok{/}\FloatTok{10}\NormalTok{) }\OperatorTok{.*} \VariableTok{exp}\NormalTok{(}\OperatorTok{{-}}\FloatTok{10}\OperatorTok{.\^{}}\NormalTok{(}\VariableTok{y}\OperatorTok{/}\FloatTok{10}\NormalTok{))}\OperatorTok{;}
\VariableTok{plot}\NormalTok{(}\VariableTok{y}\OperatorTok{,} \VariableTok{f\_Y}\OperatorTok{,} \SpecialStringTok{\textquotesingle{}r\textquotesingle{}}\OperatorTok{,} \SpecialStringTok{\textquotesingle{}LineWidth\textquotesingle{}}\OperatorTok{,} \FloatTok{2}\NormalTok{)}\OperatorTok{;}
\VariableTok{xlabel}\NormalTok{(}\SpecialStringTok{\textquotesingle{}SNR (dB)\textquotesingle{}}\NormalTok{)}\OperatorTok{;} \VariableTok{ylabel}\NormalTok{(}\SpecialStringTok{\textquotesingle{}PDF\textquotesingle{}}\NormalTok{)}\OperatorTok{;}
\VariableTok{title}\NormalTok{(}\SpecialStringTok{\textquotesingle{}Distribution of SNR in dB (Rayleigh Fading)\textquotesingle{}}\NormalTok{)}\OperatorTok{;}
\VariableTok{legend}\NormalTok{(}\SpecialStringTok{\textquotesingle{}Simulation\textquotesingle{}}\OperatorTok{,} \SpecialStringTok{\textquotesingle{}Theory\textquotesingle{}}\NormalTok{)}\OperatorTok{;}
\end{Highlighting}
\end{Shaded}

\end{PSbox}

\begin{PSbox}[unbreakable]{ex}{Example 4: Rayleigh Envelope from Gaussian Components}

\tcbset{PScodemode/.style={unbreakable}}

\begin{Shaded}
\begin{Highlighting}[]
\CommentTok{\%\% Envelope of Complex Gaussian}
\VariableTok{sigma} \OperatorTok{=} \FloatTok{1}\OperatorTok{;} \VariableTok{N} \OperatorTok{=} \FloatTok{100000}\OperatorTok{;}
\VariableTok{I} \OperatorTok{=} \VariableTok{sigma} \OperatorTok{*} \VariableTok{randn}\NormalTok{(}\FloatTok{1}\OperatorTok{,} \VariableTok{N}\NormalTok{)}\OperatorTok{;}
\VariableTok{Q} \OperatorTok{=} \VariableTok{sigma} \OperatorTok{*} \VariableTok{randn}\NormalTok{(}\FloatTok{1}\OperatorTok{,} \VariableTok{N}\NormalTok{)}\OperatorTok{;}
\VariableTok{R} \OperatorTok{=} \VariableTok{sqrt}\NormalTok{(}\VariableTok{I}\OperatorTok{.\^{}}\FloatTok{2} \OperatorTok{+} \VariableTok{Q}\OperatorTok{.\^{}}\FloatTok{2}\NormalTok{)}\OperatorTok{;}

\VariableTok{figure}\OperatorTok{;}
\VariableTok{histogram}\NormalTok{(}\VariableTok{R}\OperatorTok{,} \FloatTok{100}\OperatorTok{,} \SpecialStringTok{\textquotesingle{}Normalization\textquotesingle{}}\OperatorTok{,} \SpecialStringTok{\textquotesingle{}pdf\textquotesingle{}}\NormalTok{)}\OperatorTok{;}
\VariableTok{hold} \VariableTok{on}\OperatorTok{;}
\VariableTok{r} \OperatorTok{=} \VariableTok{linspace}\NormalTok{(}\FloatTok{0}\OperatorTok{,} \FloatTok{5}\OperatorTok{,} \FloatTok{200}\NormalTok{)}\OperatorTok{;}
\VariableTok{plot}\NormalTok{(}\VariableTok{r}\OperatorTok{,} \VariableTok{raylpdf}\NormalTok{(}\VariableTok{r}\OperatorTok{,} \VariableTok{sigma}\NormalTok{)}\OperatorTok{,} \SpecialStringTok{\textquotesingle{}r\textquotesingle{}}\OperatorTok{,} \SpecialStringTok{\textquotesingle{}LineWidth\textquotesingle{}}\OperatorTok{,} \FloatTok{2}\NormalTok{)}\OperatorTok{;}
\VariableTok{xlabel}\NormalTok{(}\SpecialStringTok{\textquotesingle{}Envelope |h|\textquotesingle{}}\NormalTok{)}\OperatorTok{;} \VariableTok{ylabel}\NormalTok{(}\SpecialStringTok{\textquotesingle{}PDF\textquotesingle{}}\NormalTok{)}\OperatorTok{;}
\VariableTok{title}\NormalTok{(}\SpecialStringTok{\textquotesingle{}Rayleigh Distribution from I/Q Components\textquotesingle{}}\NormalTok{)}\OperatorTok{;}
\VariableTok{legend}\NormalTok{(}\SpecialStringTok{\textquotesingle{}Simulation\textquotesingle{}}\OperatorTok{,} \SpecialStringTok{\textquotesingle{}Theory\textquotesingle{}}\NormalTok{)}\OperatorTok{;}
\VariableTok{fprintf}\NormalTok{(}\SpecialStringTok{\textquotesingle{}Mean: Sim=\%.3f, Theory=\%.3f\textbackslash{}n\textquotesingle{}}\OperatorTok{,} \VariableTok{mean}\NormalTok{(}\VariableTok{R}\NormalTok{)}\OperatorTok{,} \VariableTok{sigma}\OperatorTok{*}\VariableTok{sqrt}\NormalTok{(}\VariableTok{pi}\OperatorTok{/}\FloatTok{2}\NormalTok{))}\OperatorTok{;}
\end{Highlighting}
\end{Shaded}

\end{PSbox}

\begin{PSbox}[unbreakable]{ex}{Example 5: Log-Normal from Gaussian}

\tcbset{PScodemode/.style={unbreakable}}

\begin{Shaded}
\begin{Highlighting}[]
\CommentTok{\%\% Log{-}Normal: Y = exp(X) where X \textasciitilde{} N(mu, sigma\^{}2)}
\VariableTok{mu} \OperatorTok{=} \FloatTok{0}\OperatorTok{;} \VariableTok{sigma\_x} \OperatorTok{=} \FloatTok{0.5}\OperatorTok{;} \VariableTok{N} \OperatorTok{=} \FloatTok{100000}\OperatorTok{;}
\VariableTok{X} \OperatorTok{=} \VariableTok{mu} \OperatorTok{+} \VariableTok{sigma\_x}\OperatorTok{*}\VariableTok{randn}\NormalTok{(}\FloatTok{1}\OperatorTok{,} \VariableTok{N}\NormalTok{)}\OperatorTok{;}
\VariableTok{Y} \OperatorTok{=} \VariableTok{exp}\NormalTok{(}\VariableTok{X}\NormalTok{)}\OperatorTok{;}

\VariableTok{figure}\OperatorTok{;}
\VariableTok{histogram}\NormalTok{(}\VariableTok{Y}\OperatorTok{,} \FloatTok{150}\OperatorTok{,} \SpecialStringTok{\textquotesingle{}Normalization\textquotesingle{}}\OperatorTok{,} \SpecialStringTok{\textquotesingle{}pdf\textquotesingle{}}\OperatorTok{,} \SpecialStringTok{\textquotesingle{}BinLimits\textquotesingle{}}\OperatorTok{,}\NormalTok{ [}\FloatTok{0}\OperatorTok{,} \FloatTok{5}\NormalTok{])}\OperatorTok{;}
\VariableTok{hold} \VariableTok{on}\OperatorTok{;}
\VariableTok{y} \OperatorTok{=} \VariableTok{linspace}\NormalTok{(}\FloatTok{0.01}\OperatorTok{,} \FloatTok{5}\OperatorTok{,} \FloatTok{300}\NormalTok{)}\OperatorTok{;}
\VariableTok{f\_lognorm} \OperatorTok{=} \VariableTok{lognpdf}\NormalTok{(}\VariableTok{y}\OperatorTok{,} \VariableTok{mu}\OperatorTok{,} \VariableTok{sigma\_x}\NormalTok{)}\OperatorTok{;}
\VariableTok{plot}\NormalTok{(}\VariableTok{y}\OperatorTok{,} \VariableTok{f\_lognorm}\OperatorTok{,} \SpecialStringTok{\textquotesingle{}r\textquotesingle{}}\OperatorTok{,} \SpecialStringTok{\textquotesingle{}LineWidth\textquotesingle{}}\OperatorTok{,} \FloatTok{2}\NormalTok{)}\OperatorTok{;}
\VariableTok{xlabel}\NormalTok{(}\SpecialStringTok{\textquotesingle{}Y = e\^{}X\textquotesingle{}}\NormalTok{)}\OperatorTok{;} \VariableTok{ylabel}\NormalTok{(}\SpecialStringTok{\textquotesingle{}PDF\textquotesingle{}}\NormalTok{)}\OperatorTok{;}
\VariableTok{title}\NormalTok{(}\SpecialStringTok{\textquotesingle{}Log{-}Normal Distribution\textquotesingle{}}\NormalTok{)}\OperatorTok{;}
\VariableTok{legend}\NormalTok{(}\SpecialStringTok{\textquotesingle{}Simulation\textquotesingle{}}\OperatorTok{,} \SpecialStringTok{\textquotesingle{}Theory\textquotesingle{}}\NormalTok{)}\OperatorTok{;}
\end{Highlighting}
\end{Shaded}

\end{PSbox}

\PSsection{5.8}{Practice Problems}

\begin{PSbox}[unbreakable]{prob}{Problem 1}

Signal X \textasciitilde{} Uniform[0, 1]. Output Y = -2\ensuremath{\cdot}ln(X). Find the PDF of Y and identify the distribution.

\PSsolution{} CDF method: F\textsubscript{Y}(y) = P(-2ln(X) \ensuremath{\le} y) = P(X \ensuremath{\ge} e\textsuperscript{-y/2}) = 1 - e\textsuperscript{-y/2} for y \ensuremath{\ge} 0.\PSbr{}f\textsubscript{Y}(y) = (1/2)e\textsuperscript{-y/2} \ensuremath{\to} Y \textasciitilde{} Exponential(\ensuremath{\beta} = 2). (This is the inverse CDF method for generating exponentials!)

\end{PSbox}

\begin{PSbox}[unbreakable]{prob}{Problem 2}

Noise voltage X \textasciitilde{} N(0, \ensuremath{\sigma}\textsuperscript{2}) with \ensuremath{\sigma} = 2V. Power dissipated in R = 50\ensuremath{\Omega}: P = X\textsuperscript{2}/R.

\begin{enumerate}
\def\labelenumi{(\alph{enumi})}
\tightlist
\item
  Find E[P]. (b) Find the PDF of P. (c) Find P(P \textgreater{} 0.2W).
\end{enumerate}

\PSsolution{} 

\begin{enumerate}
\def\labelenumi{(\alph{enumi})}
\tightlist
\item
  E[P] = E[X\textsuperscript{2}]/R = \ensuremath{\sigma}\textsuperscript{2}/R = 4/50 = 0.08W
\item
  P = X\textsuperscript{2}/50. Let W = X\textsuperscript{2} \textasciitilde{} \ensuremath{\sigma}\textsuperscript{2}\ensuremath{\cdot}\ensuremath{\chi}\textsuperscript{2}(1). Then P = W/50. f\textsubscript{P}(p) = 50\ensuremath{\cdot}f\textsubscript{W}(50p) = (50/(2\ensuremath{\sigma}\textsuperscript{2}))\ensuremath{\cdot}(50p/\ensuremath{\sigma}\textsuperscript{2})\textsuperscript{(-1/2)\ensuremath{\cdot}e}(-50p/(2\ensuremath{\sigma}\textsuperscript{2})) for p \textgreater{} 0.
\item
  Use MATLAB: \PScode{1 - chi2cdf(0.2*50/4, 1)} = 1 - chi2cdf(2.5, 1) \ensuremath{\approx} 0.114
\end{enumerate}

\end{PSbox}

\begin{PSbox}[unbreakable]{prob}{Problem 3}

X \textasciitilde{} Exponential(\ensuremath{\lambda} = 1). Y = \ensuremath{\surd}X. Find f\textsubscript{Y}(y).

\PSsolution{} g(x) = \ensuremath{\surd}x \ensuremath{\to} x = y\textsuperscript{2}, dx/dy = 2y.\PSbr{}f\textsubscript{Y}(y) = f\textsubscript{X}(y\textsuperscript{2})\ensuremath{\cdot}\textbar{}2y\textbar{} = e\textsuperscript{-y\textsuperscript{2}}\ensuremath{\cdot}2y for y \ensuremath{\ge} 0. This is a Rayleigh distribution with \ensuremath{\sigma}\textsuperscript{2} = 1/2.

\end{PSbox}

\begin{PSbox}[unbreakable]{prob}{Problem 4}

Distance D \textasciitilde{} Uniform[1, 10] km. Path loss: L = 20\ensuremath{\cdot}log\textsubscript{10}(D) dB. Find E[L] and the PDF of L.

\PSsolution{} E[L] = E[20\ensuremath{\cdot}log\textsubscript{10}(D)] = \ensuremath{\int}\textsubscript{1}\textsuperscript{10} 20\ensuremath{\cdot}log\textsubscript{10}(d)\ensuremath{\cdot}(1/9) dd = (20/9)\ensuremath{\cdot}\ensuremath{\int}\textsubscript{1}\textsuperscript{10} log\textsubscript{10}(d) dd\PSbr{}= (20/9)\ensuremath{\cdot}[d\ensuremath{\cdot}log\textsubscript{10}(d) - d/ln(10)]\textsubscript{1}\textsuperscript{10} = (20/9)\ensuremath{\cdot}[10 - 10/ln10 + 1/ln10] = (20/9)\ensuremath{\cdot}(10 - 9/ln10) \ensuremath{\approx} 13.55 dB

PDF: L ranges from 0 to 20 dB. g\textsuperscript{-1}(l) = 10\textsuperscript{l/20}, \textbar{}dg\textsuperscript{-1}/dl\textbar{} = (ln10/20)\ensuremath{\cdot}10\textsuperscript{l/20}.\PSbr{}f\textsubscript{L}(l) = (1/9)\ensuremath{\cdot}(ln10/20)\ensuremath{\cdot}10\textsuperscript{l/20} for 0 \ensuremath{\le} l \ensuremath{\le} 20.

\end{PSbox}

\begin{PSbox}[unbreakable]{prob}{Problem 5}

Write MATLAB code to verify Problem 2 by simulation.

\PSsolution{} 

\tcbset{PScodemode/.style={unbreakable}}

\begin{Shaded}
\begin{Highlighting}[]
\VariableTok{sigma} \OperatorTok{=} \FloatTok{2}\OperatorTok{;} \VariableTok{R} \OperatorTok{=} \FloatTok{50}\OperatorTok{;} \VariableTok{N} \OperatorTok{=} \FloatTok{500000}\OperatorTok{;}
\VariableTok{X} \OperatorTok{=} \VariableTok{sigma}\OperatorTok{*}\VariableTok{randn}\NormalTok{(}\FloatTok{1}\OperatorTok{,}\VariableTok{N}\NormalTok{)}\OperatorTok{;}
\VariableTok{P} \OperatorTok{=} \VariableTok{X}\OperatorTok{.\^{}}\FloatTok{2} \OperatorTok{/} \VariableTok{R}\OperatorTok{;}
\VariableTok{fprintf}\NormalTok{(}\SpecialStringTok{\textquotesingle{}E[P]: Theory=\%.4f W, Sim=\%.4f W\textbackslash{}n\textquotesingle{}}\OperatorTok{,} \VariableTok{sigma}\OperatorTok{\^{}}\FloatTok{2}\OperatorTok{/}\VariableTok{R}\OperatorTok{,} \VariableTok{mean}\NormalTok{(}\VariableTok{P}\NormalTok{))}\OperatorTok{;}
\VariableTok{fprintf}\NormalTok{(}\SpecialStringTok{\textquotesingle{}P(P\textgreater{}0.2): Theory=\%.4f, Sim=\%.4f\textbackslash{}n\textquotesingle{}}\OperatorTok{,} \FloatTok{1}\OperatorTok{{-}}\VariableTok{chi2cdf}\NormalTok{(}\FloatTok{0.2}\OperatorTok{*}\VariableTok{R}\OperatorTok{/}\VariableTok{sigma}\OperatorTok{\^{}}\FloatTok{2}\OperatorTok{,}\FloatTok{1}\NormalTok{)}\OperatorTok{,} \VariableTok{mean}\NormalTok{(}\VariableTok{P}\OperatorTok{\textgreater{}}\FloatTok{0.2}\NormalTok{))}\OperatorTok{;}
\end{Highlighting}
\end{Shaded}

\emph{Next Module: {Module 6 — Two Random Variables}}

\end{PSbox}
